\section{Apéndice de datos}\label{app:data}

\subsection{Muestra y construcción de variables}

Agrupo los veinte módulos trimestrales de empleo de la ENAHO (módulo 500) para 2021Q1--2025Q4; el análisis principal usa 2021Q1--2024Q4 y la validación usa 2024Q1--2025Q4. La muestra en edad de trabajar comprende a las personas de 14 a 65 años con información de educación. Todas las estadísticas usan el factor de expansión de empleo del módulo (\texttt{fac500a}). Este factor está calibrado para que el año completo de entrevistas se expanda a la proyección anual de población; cada submuestra trimestral se expande, por tanto, a aproximadamente la cuarta parte de la población, lo que es irrelevante para las medias, proporciones y estimaciones de regresión ponderadas que uso a lo largo del trabajo (numerador y denominador se escalan juntos) y solo importaría para totales poblacionales, que nunca estimo. La partición trimestral sigue el mes de la entrevista y coincide con el diseño de la propia encuesta: la inferencia nacional trimestral es un nivel de inferencia oficial de la ENAHO en las fichas técnicas tanto de 2021 como de 2025, y la muestra no panel vuelve a visitar los mismos conglomerados en los mismos meses cada año. Verifico ambas propiedades en los microdatos para toda la ventana: los 5359 conglomerados aparecen en todos los años de 2021 a 2025, la superposición de unidades primarias de muestreo entre el mismo trimestre de años consecutivos va de 70 a 92 por ciento (frente a 2 a 5 por ciento entre trimestres distintos), y cada trimestre es una réplica internamente balanceada del diseño. Conviene registrar dos cambios en la encuesta durante la ventana. Primero, la modalidad mixta de entrevista adoptada durante la pandemia (prioridad a la entrevista presencial, con seguimiento telefónico para los miembros ausentes) siguió siendo el estándar hasta 2025, de modo que la modalidad es en líneas generales consistente en mi periodo de estudio; la única excepción es febrero de 2021, cuando un cuestionario telefónico reducido en quince regiones omitió, entre otras preguntas, la del tipo de contrato (faltante para la mitad de los asalariados de ese mes), lo que no afecta a ninguno de los resultados de base (la informalidad de 2021 proviene del indicador oficial del INEI y las preguntas de ingreso se mantuvieron) y se trata explícitamente en el único margen del apéndice que usa contratos. Segundo, la ficha técnica de 2025 describe el diseño como bietápico, con conglomerados urbanos de unas 140 viviendas, frente a la descripción polietápica de 2021; los identificadores y conteos de conglomerados son, no obstante, continuos a lo largo de la ventana, de modo que se trata de una reformulación de la documentación y no de un quiebre en la muestra. Los departamentos son las 25 unidades administrativas de primer nivel derivadas de los dos primeros dígitos del \texttt{ubigeo}; el indicador urbano se deriva del estrato de muestreo.

\paragraph{Educación y grupos de calificación.} El nivel educativo alcanzado corresponde a la variable \texttt{p301a} de la ENAHO. Verifiqué su codificación con la distribución ponderada del nivel alcanzado. Los trabajadores poco calificados son aquellos con \texttt{p301a} entre 1 y 6 (desde sin nivel hasta secundaria completa); los trabajadores calificados son aquellos con \texttt{p301a} entre 7 y 11 (desde superior no universitaria hasta posgrado). Se excluyen la educación básica especial (código 12) y los valores faltantes.

\paragraph{Ingreso laboral y salarios.} La ENAHO registra la remuneración de los asalariados por periodo de pago: \texttt{p524a1} es el pago bruto, incluidas las horas extra, y \texttt{p523} su frecuencia (diaria, semanal, quincenal o mensual). El cincuenta y cuatro por ciento de los asalariados de baja educación recibe pagos con frecuencia no mensual, frente al quince por ciento de los asalariados de alta educación, por lo que convierto todos los pagos a equivalentes mensuales con los mismos factores implícitos en la variable de ingreso anualizado del propio INEI (\texttt{i524a1}): 25 para la frecuencia diaria, $52/12$ para la semanal, 2 para la quincenal y 1 para la mensual. La distribución ponderada completa de frecuencias entre los asalariados privados con ingresos válidos es de 2 por ciento diaria, 45 semanal, 9 quincenal y 45 mensual para el grupo de baja educación, frente a 1, 18, 6 y 76 por ciento para el grupo de alta educación; la variable de frecuencia nunca falta en esa muestra. Los asalariados comprenden a empleados y trabajadores del hogar (\texttt{p507} en $\{3,4,6\}$); los trabajadores del hogar tienen derecho al salario mínimo desde la Ley 31047. Para los empleadores y los trabajadores por cuenta propia, el ingreso laboral es la ganancia neta (\texttt{p530a}), que ya es mensual. Como el ingreso reportado se refiere al mes anterior a la entrevista, deflacto en ese mes de referencia (las entrevistas de enero se remontan a diciembre del año anterior) con el índice de precios al consumidor de Lima con base diciembre de 2021, y evalúo el mínimo legal vigente en ese mismo mes de referencia al clasificar a los trabajadores como remunerados por debajo del piso. El salario real por hora divide el ingreso laboral mensual equivalente entre $4.345$ veces las horas semanales en la ocupación principal (\texttt{p513t}; la variable de horas imputadas del INEI, \texttt{i513t}, la reemplaza cuando \texttt{p513t} falta, pero ninguna observación de la muestra de análisis salarial y menos del 0.01 por ciento de la muestra de horas lo requieren). Los resultados salariales excluyen a los trabajadores del sector público, cuya remuneración se fija mediante escalas propias, y se winsorizan una sola vez en los percentiles 0.5 y 99.5 de la muestra agrupada de asalariados privados.

\paragraph{Informalidad.} El indicador oficial de empleo informal del INEI (\texttt{ocupinf}, codificado como 1 para informal y 2 para formal en estos archivos) está disponible para 2021--2023 pero no para 2024--2025. Reconstruyo un indicador consistente a partir de variables primitivas disponibles en todos los años. Un empleado asalariado o trabajador del hogar se clasifica como formal si está afiliado a un sistema de pensiones (\texttt{p558a1}--\texttt{p558a4}, cuyas opciones seleccionadas se marcan con su código de opción y no con unos) o si tiene un contrato laboral formal (\texttt{p511a} entre 1 y 6); un empleador o trabajador por cuenta propia se clasifica como formal si la unidad productiva está registrada (\texttt{p510a1}); los trabajadores familiares no remunerados son siempre informales. Validada contra \texttt{ocupinf} en los años de superposición 2021--2023, esta regla coincide con la clasificación oficial para el 87.8 por ciento de los ocupados y arroja una tasa de informalidad de 74.9 por ciento, frente al 79.1 por ciento oficial. Uso el indicador oficial siempre que está disponible y la reconstrucción solo para 2024--2025; dentro de la ventana de validación 2024--2025 el indicador se reconstruye en todos los trimestres, de modo que allí se mide de manera consistente.

\paragraph{Salario mínimo e incidencia.} El calendario del salario mínimo legal se toma de los decretos y se verifica con la serie del Banco Central de Reserva del Perú (PN02124PM); el deflactor del IPC es la serie PN38705PM del BCRP (Lima, diciembre de 2021 $=$ 100), y ambos insumos se regeneran con un script documentado en el paquete de replicación. El índice de Kaitz departamental es el cociente entre el mínimo posterior a la reforma (1025 soles) y el salario mensual equivalente mediano ponderado de los asalariados privados del departamento antes de la reforma (2021Q1--2022Q1). La medida de fracción afectada es la proporción ponderada, antes de la reforma, de asalariados privados que ganaban menos de 1025 soles.

\paragraph{El Panel ENAHO.} El análisis de transiciones de la Sección~\ref{sec:transitions} usa el Panel ENAHO 2020--2024 (encuesta 978 del INEI, publicada en noviembre de 2025), cuyo módulo de empleo (\texttt{enaho01a-2020-2024-500-panel.dta}) contiene un registro por persona con variables con sufijo de año para 2020 a 2024. Las personas observadas en ambos años de un par se identifican con los indicadores \texttt{perpanel} (23\,983 personas enlazadas para 2021--22 y 23\,245 para 2022--23) y se ponderan con los factores de expansión de panel del INEI, \texttt{facpanel}, que corrigen por desgaste y no respuesta; la inferencia oficial para pares de años consecutivos es nacional, urbana y rural, y por región natural. Los estados de empleo replican exactamente las definiciones de corte transversal: ocupado es \texttt{ocu500}${}=1$; el estado formal privado requiere empleo en el sector privado con un puesto formal; la informalidad usa el \texttt{ocupinf} oficial donde se publica y, en los demás casos, la reconstrucción con la Regla F descrita arriba; el trabajo independiente comprende a los trabajadores por cuenta propia y a los empleadores. La educación en el año base define los grupos de calificación. Como el panel es anual y la reforma entró en vigor el 1 de mayo de 2022, las observaciones de cada par se clasifican según el mes de entrevista de 2022, como se describe en las notas del Cuadro~\ref{tab:transitions}.

\paragraph{Intensidad de las protestas.} La intensidad departamental de la crisis política de 2022--23 se mide a partir de los reportes mensuales de conflictos de la Presidencia del Consejo de Ministros (Secretar\'ia de Gesti\'on Social y Di\'alogo, ``Reporte de Conflictos Sociales''), cuyos cuadros de seguimiento de la crisis registran el número máximo de personas movilizadas en acciones de protesta en cada departamento, por semana y por mes. Tomo el máximo de cada departamento en los reportes de enero y febrero de 2023, los dos meses en que la movilización nacional alcanzó su punto más alto (el reporte de diciembre de 2022 es anterior al cuadro), lo divido entre la población en edad de trabajar del departamento (cuatro veces la suma trimestral media de los factores de expansión de la ENAHO) y estandarizo el resultado entre los 25 departamentos. El cuadro transcrito y los documentos fuente se incluyen en el paquete de replicación.

\subsection{Resultados complementarios}

El Cuadro~\ref{tab:ddd} presenta las estimaciones de triple diferencia y el Cuadro~\ref{tab:loo} los rangos de estabilidad excluyendo un departamento a la vez discutidos en la Sección~\ref{sec:robustness}; los Cuadros~\ref{tab:expocells}--\ref{tab:margins} reúnen los diagnósticos del diseño de exposición, las verificaciones sectoriales y ocupacionales, las cotas de selección y los cálculos de potencia, los resultados de diferencias en diferencias sintéticas y de potencia de la prueba previa, y los márgenes de ajuste. El Cuadro~\ref{tab:crisis} presenta la robustez del diseño de exposición regional continua frente a la crisis política y el Cuadro~\ref{tab:uncond} los resultados de composición incondicionales, ambos discutidos en la Sección~\ref{sec:robustness}. El Cuadro~\ref{tab:placebo} presenta las reformas placebo en el tiempo y confirma que los placebos salariales son nulos, mientras que los placebos de empleo están contaminados por la tendencia previa de la pandemia. Las estimaciones doblemente robustas de \citet{santanna_zhao_2020} aparecen junto a las estimaciones con efectos fijos bidireccionales en el Cuadro~\ref{tab:main}. El Figura~\ref{fig:val2025} muestra el estudio de eventos en torno a la reforma de enero de 2025 usado en la validación fuera de muestra.

\begin{table}[!tbp]\centering
\caption{Estimaciones de triple diferencia: $\mathrm{Low}\times\mathrm{Post}\times\mathrm{HighBite}$}\label{tab:ddd}
\begin{threeparttable}
\begin{tabular}{lc}
\toprule
Resultado & Estimación DDD \\
\midrule
Log salario/hora & -0.027 \\
 & (0.038) \\
Empleo & 0.015 \\
 & (0.016) \\
Empleo formal & -0.018 \\
 & (0.018) \\
Trab. independiente & 0.026 \\
 & (0.018) \\
\bottomrule
\end{tabular}
\begin{tablenotes}\footnotesize
\item Notas: Coeficiente de la interacción triple $\mathrm{Low}\times\mathrm{Post}\times\mathrm{HighBite}$, donde HighBite identifica a los departamentos cuya fracción de asalariados privados que ganaba menos que el nuevo salario mínimo antes de la reforma está por encima de la mediana. Efectos fijos de departamento y trimestre, controles como en la especificación de base; errores estándar agrupados por departamento (25 grupos); factores de expansión de la ENAHO. $^{*}p<0.1$, $^{**}p<0.05$, $^{***}p<0.01$.
\end{tablenotes}
\end{threeparttable}\end{table}


\begin{table}[!tbp]\centering
\caption{Estabilidad de las estimaciones DiD de base al excluir una unidad a la vez}\label{tab:loo}
\begin{threeparttable}
\begin{tabular}{lccc}
\toprule
Resultado & Muestra completa & Rango sin un depto. & Rango sin un trimestre \\
\midrule
Log salario/hora & -0.021 & [-0.027, -0.017] & [-0.026, -0.010] \\
Empleo & -0.025 & [-0.028, -0.018] & [-0.031, -0.014] \\
Empleo formal & -0.046 & [-0.050, -0.043] & [-0.052, -0.042] \\
Trab. independiente & 0.036 & [0.032, 0.038] & [0.031, 0.041] \\
\bottomrule
\end{tabular}
\begin{tablenotes}\footnotesize
\item Notas: Rangos del coeficiente $\mathrm{Low}\times\mathrm{Post}$ al reestimar la especificación de base excluyendo un departamento a la vez (25 estimaciones) o un trimestre a la vez (15 estimaciones). Las estimaciones completas por unidad están en los resultados de replicación (\texttt{rob\_leaveoneout\_detail.csv}).
\end{tablenotes}
\end{threeparttable}\end{table}


\begin{table}[!tbp]\centering
\caption{Diseños de exposición: diagnósticos y estimaciones a nivel de celda}\label{tab:expocells}
\begin{threeparttable}
\footnotesize
\begin{tabular}{lcccc}
\toprule
\multicolumn{5}{l}{\emph{Panel A: diseño regional de Kaitz, pruebas conjuntas de adelantos de la exposición}} \\
\midrule
Log salario/hora & \multicolumn{4}{c}{adelantos $p$ = 0.243} \\
Empleo & \multicolumn{4}{c}{adelantos $p$ = 0.002} \\
Empleo formal & \multicolumn{4}{c}{adelantos $p$ = 0.842} \\
Trab. independiente & \multicolumn{4}{c}{adelantos $p$ = 0.202} \\
\midrule
\multicolumn{5}{l}{\emph{Panel B: diseños de fracción afectada a nivel de celda (ejercicio de transparencia)}} \\
Resultado & Exposición & Estimación & (s.e.) & $p$ adelantos \\ \midrule
Log salario/hora & frac\_below & 0.129$^{**}$ & (0.056) & 0.011 \\
Log salario/hora & frac\_aff & 0.016 & (0.133) & 0.004 \\
Empleo & frac\_below & -0.225$^{***}$ & (0.032) & 0.000 \\
Empleo & frac\_aff & 0.391$^{***}$ & (0.072) & 0.000 \\
Empleo formal & frac\_below & -0.074$^{***}$ & (0.014) & 0.007 \\
Empleo formal & frac\_aff & 0.303$^{***}$ & (0.039) & 0.031 \\
Trab. independiente & frac\_below & 0.079$^{***}$ & (0.018) & 0.244 \\
Trab. independiente & frac\_aff & -0.155$^{**}$ & (0.065) & 0.045 \\
\bottomrule
\end{tabular}
\begin{tablenotes}\footnotesize
\item Notas: Panel A: pruebas conjuntas de que las interacciones trimestre previo a la reforma$\times$exposición son cero en el diseño regional de exposición continua (índice de Kaitz departamental estandarizado). Panel B: coeficientes de $\mathrm{Post}\times$exposición en diseños de fracción afectada a nivel de celda (celdas departamento$\times$educación$\times$edad; frac\_below = proporción de los asalariados privados de la celda que ganaba menos de 1025 antes de la reforma; frac\_aff = proporción que ganaba entre 930 y 1025), con efectos fijos de celda y trimestre y errores estándar agrupados por departamento. Ambas variantes a nivel de celda fallan sus diagnósticos de adelantos (la exposición de grano fino se confunde con las trayectorias de recuperación pospandemia de las celdas), por lo que el diseño regional es la corroboración preferida. $^{*}p<0.1$, $^{**}p<0.05$, $^{***}p<0.01$.
\end{tablenotes}
\end{threeparttable}\end{table}


\begin{table}[!tbp]\centering
\caption{Amplitud sectorial y dosis-respuesta ocupacional de la reasignación}\label{tab:sectorocc}
\begin{threeparttable}
\footnotesize
\begin{tabular}{lccc}
\toprule
\multicolumn{4}{l}{\emph{Panel A: DiD por educación según grupo sectorial}} \\
Resultado & Cubiertos de alta incidencia & Agricultura & Otros privados \\
\midrule
Empleo formal & -0.029$^{***}$ & -0.051$^{***}$ & -0.044$^{***}$ \\
 & (0.006) & (0.013) & (0.012) \\
Trab. independiente & 0.028$^{***}$ & 0.020 & 0.030$^{**}$ \\
 & (0.007) & (0.013) & (0.012) \\
Log salario/hora & -0.040$^{**}$ & -0.060 & 0.015 \\
 & (0.018) & (0.050) & (0.014) \\
\midrule
\multicolumn{4}{l}{\emph{Panel B: dosis-respuesta según exposición ocupacional ($\mathrm{Post}\times$ proporción de la ocupación bajo 1025)}} \\
Resultado & Estimación & (s.e.) & $p$ adelantos \\ \midrule
Empleo formal & -0.049$^{***}$ & (0.016) & 0.008 \\
Trab. independiente & 0.000 & (0.017) & 0.036 \\
Log salario/hora & 0.131$^{***}$ & (0.022) & 0.106 \\
Gana menos que el SM & -0.094$^{***}$ & (0.023) & 0.000 \\
\bottomrule
\end{tabular}
\begin{tablenotes}\footnotesize
\item Notas: Panel A: DiD de $\mathrm{Low}\times\mathrm{Post}$ dentro de grupos sectoriales definidos a partir de las divisiones de la CIIU Rev.4 (cubiertos de alta incidencia: manufactura 10--33, construcción 41--43, comercio 45--47, hoteles y restaurantes 55--56; agricultura: divisiones 01--03, bajo el régimen agrario de la Ley 31110, cuya remuneración diaria está indexada a la RMV y, por tanto, subió con la reforma). Panel B: coeficiente de $\mathrm{Post}\times$(proporción, antes de la reforma, de los asalariados privados de la ocupación a 2 dígitos que ganaba menos de 1025), con efectos fijos de ocupación, trimestre y departamento; la ocupación se mide en la entrevista y, por tanto, es posterior al tratamiento para quienes cambian de ocupación, de modo que el panel B es una corroboración descriptiva y no un diseño principal. Errores estándar agrupados por departamento; factores de expansión de la ENAHO. $^{*}p<0.1$, $^{**}p<0.05$, $^{***}p<0.01$.
\end{tablenotes}
\end{threeparttable}\end{table}


\begin{table}[!tbp]\centering
\caption{Cotas por selección y poder estadístico}\label{tab:leemde}
\begin{threeparttable}
\footnotesize
\begin{tabular}{lcc}
\toprule
\multicolumn{3}{l}{\emph{Panel A: cotas de recorte tipo Lee, DiD 2$\times$2 incondicional del log del salario}} \\
\midrule
DiD de selección (inclusión en la muestra salarial) & -0.021 & \\
Fracción de recorte & 0.076 & \\
DiD sin recorte & -0.005 & \\
Cotas & [-0.099, 0.101] & \\
\midrule
\multicolumn{3}{l}{\emph{Panel B: efectos mínimos detectables (poder de 80\%, tamaño de 5\%)}} \\
Resultado & Estimación & EMD \\ \midrule
Log salario real por hora & -0.021 & 0.034 \\
Log ingreso real mensual & -0.010 & 0.038 \\
Ocupado (edad de trabajar) & -0.025 & 0.016 \\
Participación laboral & -0.008 & 0.012 \\
Empleo formal & -0.046 & 0.021 \\
Horas semanales (empleo principal) & 0.556 & 0.698 \\
Tiempo completo (>=35h) & -0.013 & 0.013 \\
Independiente & 0.036 & 0.017 \\
Gana menos que el SM & -0.006 & 0.021 \\
\bottomrule
\end{tabular}
\begin{tablenotes}\footnotesize
\item Notas: Panel A: cotas para el DiD 2$\times$2 incondicional del log del salario por hora, siguiendo la lógica de recorte de \citet{lee_2009}: la reforma redujo la probabilidad de que un trabajador de baja educación aparezca en la muestra de asalariados (DiD de selección), por lo que la celda low$\times$post se recorta por arriba y por abajo en la fracción implícita de su tasa base. Panel B: efecto mínimo detectable $=2.8\times$ el error estándar agrupado por departamento del DiD de base.
\end{tablenotes}
\end{threeparttable}\end{table}


\begin{table}[!tbp]\centering
\caption{Diferencias en diferencias sintéticas y poder de la prueba de tendencias previas}\label{tab:sdidpower}
\begin{threeparttable}
\footnotesize
\begin{tabular}{lccc}
\toprule
\multicolumn{4}{l}{\emph{Panel A: DiD sintéticas (deptos. de alta incidencia vs.\ sintéticos de baja)}} \\
Resultado & Medias, baja educ. & Brecha educativa & \\
\midrule
Log salario/hora & 0.041 & -0.028 & \\
 & (0.025) & (0.048) & \\
Empleo & -0.039$^{***}$ & -0.016 & \\
 & (0.010) & (0.011) & \\
Empleo formal & -0.013$^{**}$ & -0.028 & \\
 & (0.007) & (0.017) & \\
Trab. independiente & 0.011 & 0.023 & \\
 & (0.011) & (0.019) & \\
\midrule
\multicolumn{4}{l}{\emph{Panel B: poder de la prueba de tendencias previas frente a violaciones lineales \citep{roth_2022_pretest}}} \\
Resultado & Pendiente (50\% potencia) & Sesgo impl. (prom.\ post) & Sesgo / estimación \\ \midrule
Log salario/hora & 0.008 & 0.052 & -2.53 \\
Empleo & 0.006 & 0.036 & -1.43 \\
Empleo formal & 0.006 & 0.041 & -0.89 \\
Trab. independiente & 0.006 & 0.041 & 1.13 \\
\bottomrule
\end{tabular}
\begin{tablenotes}\footnotesize
\item Notas: Panel A: estimaciones de diferencias en diferencias sintéticas \citep{arkhangelsky_2021_sdid} sobre celdas departamento$\times$trimestre, que tratan como tratados desde 2022Q3 a los departamentos con incidencia por encima de la mediana; la columna (1) usa las medias de los resultados de baja educación y la columna (2), la brecha entre baja y alta educación; errores estándar jackknife entre paréntesis (la varianza placebo no es factible con 12 unidades tratadas y 13 de control). Panel B: pendiente de una violación lineal de las tendencias paralelas frente a la cual la prueba previa convencional tiene solo 50 por ciento de poder \citep{roth_2022_pretest}, el sesgo que tal violación no detectada induciría en el efecto promedio posterior y su razón respecto del efecto estimado. $^{*}p<0.1$, $^{**}p<0.05$, $^{***}p<0.01$.
\end{tablenotes}
\end{threeparttable}\end{table}


\begin{table}[!tbp]\centering
\caption{Márgenes de ajuste y subgrupos adicionales}\label{tab:margins}
\begin{threeparttable}
\footnotesize
\begin{tabular}{lccc}
\toprule
\multicolumn{4}{l}{\emph{Panel A: márgenes de ajuste ($\mathrm{Low}\times\mathrm{Post}$)}} \\
Margen & Estimación & (s.e.) & Media previa \\
\midrule
Sin contrato escrito (asalariados) & 0.031$^{***}$ & (0.008) & 0.610 \\
Microempresa, <=20 trab. (ocupados) & 0.028$^{***}$ & (0.007) & 0.854 \\
Antigüedad < 12 meses (asalariados) & -0.012 & (0.016) & 0.513 \\
Informal en sector informal (2021-23) & 0.024$^{**}$ & (0.010) & 0.663 \\
Informal en sector FORMAL (2021-23) & 0.008 & (0.007) & 0.144 \\
\midrule
\multicolumn{4}{l}{\emph{Panel B: subgrupos por rol en el hogar y autoidentificación étnica}} \\
Resultado & Grupo & Estimación & (s.e.) \\ \midrule
Empleo & Jefe de hogar & -0.003 & (0.006) \\
Empleo & No jefe de hogar & -0.038$^{***}$ & (0.008) \\
Empleo & Indígena & -0.004 & (0.015) \\
Empleo & No indígena & -0.030$^{***}$ & (0.005) \\
Empleo formal & Jefe de hogar & -0.005 & (0.011) \\
Empleo formal & No jefe de hogar & -0.068$^{***}$ & (0.006) \\
Empleo formal & Indígena & -0.030 & (0.018) \\
Empleo formal & No indígena & -0.048$^{***}$ & (0.008) \\
Trab. independiente & Jefe de hogar & 0.030$^{***}$ & (0.008) \\
Trab. independiente & No jefe de hogar & 0.039$^{***}$ & (0.008) \\
Trab. independiente & Indígena & 0.009 & (0.017) \\
Trab. independiente & No indígena & 0.045$^{***}$ & (0.008) \\
Log salario/hora & Jefe de hogar & -0.020 & (0.017) \\
Log salario/hora & No jefe de hogar & -0.019 & (0.014) \\
Log salario/hora & Indígena & -0.008 & (0.029) \\
Log salario/hora & No indígena & -0.026$^{*}$ & (0.013) \\
\bottomrule
\end{tabular}
\begin{tablenotes}\footnotesize
\item Notas: Panel A: estimaciones DiD de $\mathrm{Low}\times\mathrm{Post}$ para resultados orientados a mecanismos: la proporción de asalariados privados sin contrato escrito (\texttt{p511a}), la proporción de trabajadores privados en unidades de a lo sumo 20 trabajadores (\texttt{p512a}), la proporción de asalariados con menos de 12 meses de antigüedad (\texttt{p513a1/2}) y la descomposición del empleo informal en empleos informales del sector informal y del sector formal (INEI \texttt{emplpsec}, disponible para 2021--2023). Panel B: DiD de base dentro de subgrupos definidos por la jefatura del hogar (\texttt{p203}) y la autoidentificación indígena (\texttt{p558c}). Errores estándar agrupados por departamento; factores de expansión de la ENAHO. $^{*}p<0.1$, $^{**}p<0.05$, $^{***}p<0.01$.
\end{tablenotes}
\end{threeparttable}\end{table}


\begin{table}[!tbp]\centering
\caption{Robustez frente a la ola migratoria venezolana}\label{tab:migration}
\begin{threeparttable}
\footnotesize
\setlength{\tabcolsep}{4pt}
\begin{tabular}{lcccc}
\toprule
Especificación & Log salario/hora & Empleo & Empleo formal & Trab. independiente \\
\midrule
Base & -0.021 & -0.025$^{***}$ & -0.046$^{***}$ & 0.036$^{***}$ \\
 & (0.012) & (0.006) & (0.007) & (0.006) \\
EF departamento$\times$trimestre & -0.020 & -0.015$^{**}$ & -0.041$^{***}$ & 0.034$^{***}$ \\
 & (0.013) & (0.006) & (0.008) & (0.006) \\
Excluyendo Lima y Callao & -0.021 & -0.018$^{*}$ & -0.042$^{***}$ & 0.037$^{***}$ \\
 & (0.022) & (0.009) & (0.010) & (0.010) \\
Sin los 8 deptos. con más migrantes & -0.051 & -0.012 & -0.049$^{***}$ & 0.045$^{**}$ \\
 & (0.032) & (0.013) & (0.015) & (0.016) \\
\bottomrule
\end{tabular}
\begin{tablenotes}\footnotesize
\item Notas: Estimaciones DiD de $\mathrm{Low}\times\mathrm{Post}$ bajo especificaciones que abordan la llegada de inmigrantes venezolanos. Los efectos fijos departamento$\times$trimestre absorben cualquier choque local común a ambos grupos educativos dentro de un departamento-trimestre, incluidas las llegadas de migrantes y las olas de regularización. La exclusión de los 8 principales elimina Lima, Callao, La Libertad, Arequipa, Lambayeque, Piura, Ica y Tumbes, los departamentos que acogen a la gran mayoría de la población venezolana. Errores estándar agrupados por departamento; factores de expansión de la ENAHO. $^{*}p<0.1$, $^{**}p<0.05$, $^{***}p<0.01$.
\end{tablenotes}
\end{threeparttable}\end{table}


\begin{table}[!tbp]\centering
\caption{La crisis política de 2022--23 y el diseño de exposición regional continua}\label{tab:crisis}
\begin{threeparttable}
\footnotesize
\setlength{\tabcolsep}{3pt}
\begin{tabular}{lcccccc}
\toprule
& (1) Base & (2) +Post$\times$ & (3) +2023$\times$ & (4) Excl.\ bloque & (5) Post = & (6) (4)+(5) \\
Resultado & & protesta & protesta & sur & solo 2022 & \\
\midrule
Log salario/hora & 0.017 & 0.012 & 0.015 & 0.009 & 0.019 & 0.008 \\
 & (0.012) & (0.012) & (0.012) & (0.012) & (0.016) & (0.014) \\
 & [0.034] & [0.074] & [0.038] & [0.156] & [0.088] & [0.355] \\
Empleo & -0.034$^{***}$ & -0.027$^{***}$ & -0.029$^{***}$ & -0.029$^{**}$ & -0.029$^{***}$ & -0.024$^{**}$ \\
 & (0.012) & (0.009) & (0.009) & (0.010) & (0.009) & (0.009) \\
 & [0.001] & [0.004] & [0.001] & [0.000] & [0.002] & [0.007] \\
Empleo formal & -0.018$^{***}$ & -0.017$^{***}$ & -0.016$^{***}$ & -0.016$^{***}$ & -0.006 & -0.003 \\
 & (0.003) & (0.003) & (0.002) & (0.003) & (0.004) & (0.003) \\
 & [0.001] & [0.001] & [0.000] & [0.002] & [0.453] & [0.877] \\
Trab. independiente & 0.010$^{**}$ & 0.008$^{**}$ & 0.008$^{**}$ & 0.010$^{**}$ & 0.007$^{*}$ & 0.008 \\
 & (0.004) & (0.004) & (0.003) & (0.004) & (0.004) & (0.006) \\
 & [0.023] & [0.052] & [0.044] & [0.060] & [0.069] & [0.151] \\
\bottomrule
\end{tabular}
\begin{tablenotes}\footnotesize
\item Notas: Cada celda reporta el coeficiente de $\mathrm{Post}\times$(índice de Kaitz departamental estandarizado) del diseño de exposición continua, con errores estándar agrupados por departamento entre paréntesis y, entre corchetes, el valor $p$ de inferencia por aleatorización (permutando los índices de Kaitz entre los departamentos incluidos, 1000 extracciones, con el control de protestas fijo en su asignación real). La columna (1) es la especificación de base del Cuadro~\ref{tab:inference}, columna (5). Las columnas (2) y (3) añaden la movilización máxima de protesta per cápita estandarizada, interactuada con el indicador posterior a la reforma y con un indicador desde 2023Q1, respectivamente; la intensidad de las protestas es el número máximo de personas movilizadas en cada departamento en enero--febrero de 2023, según los reportes mensuales de conflictos de la Presidencia del Consejo de Ministros, dividido entre la población en edad de trabajar del departamento. La columna (4) excluye los seis departamentos del sur que fueron el centro de las protestas de diciembre de 2022 a marzo de 2023 (Apur\'imac, Arequipa, Ayacucho, Cusco, Madre de Dios, Puno). La columna (5) restringe la ventana posterior a la reforma a 2022Q3--Q4, antes de la destitución del presidente Castillo el 7 de diciembre de 2022, y la columna (6) combina ambas restricciones. Todas las estimaciones están ponderadas con los factores de expansión de la ENAHO. $^{*}p<0.1$, $^{**}p<0.05$, $^{***}p<0.01$.
\end{tablenotes}
\end{threeparttable}\end{table}


\begin{table}[!tbp]\centering
\caption{Resultados composicionales incondicionales (proporciones de la población en edad de trabajar)}\label{tab:uncond}
\begin{threeparttable}
\footnotesize
\setlength{\tabcolsep}{4pt}
\begin{tabular}{lccccc}
\toprule
& Media previa & DiD por educación & Tend. previa & Pendiente de Kaitz & RI \\
Resultado & a la reforma & (Low$\times$Post) & $p$ & (Post$\times z$) & $p$ \\
\midrule
Empleo formal privado & 0.148 & -0.033$^{***}$ & 0.014 & -0.016$^{***}$ & 0.000 \\
 & & (0.006) & & (0.003) & \\
Trab. independiente & 0.252 & 0.017$^{***}$ & 0.000 & -0.004$^{*}$ & 0.115 \\
 & & (0.004) & & (0.003) & \\
Empleo informal privado & 0.500 & 0.009 & 0.306 & -0.019$^{**}$ & 0.009 \\
 & & (0.006) & & (0.008) & \\
\bottomrule
\end{tabular}
\begin{tablenotes}\footnotesize
\item Notas: Los resultados se definen para toda persona en edad de trabajar: una persona cuenta como ocupada formal (privada) solo si está ocupada, en el sector privado y es formal, y de manera análoga para el trabajo independiente y el empleo informal privado; los no ocupados y los trabajadores del sector público cuentan como ceros. La columna (2) reporta el DiD por educación de base $\mathrm{Low}\times\mathrm{Post}$ con errores estándar agrupados por departamento; la columna (3), el valor $p$ conjunto de las interacciones del estudio de eventos previas a la reforma; la columna (4), el coeficiente de $\mathrm{Post}\times$(índice de Kaitz departamental estandarizado) del diseño de exposición regional continua; la columna (5), su valor $p$ de inferencia por aleatorización (permutando los índices de Kaitz entre departamentos, 1000 extracciones). Como estos resultados incorporan el margen de empleo, las versiones DiD por educación heredan su tendencia previa diferencial asociada a la recuperación pospandemia y se reportan por transparencia, sin apoyarse en ellas; las estimaciones del diseño regional no dependen del contraste por educación. Todas las estimaciones están ponderadas con los factores de expansión de la ENAHO. $^{*}p<0.1$, $^{**}p<0.05$, $^{***}p<0.01$.
\end{tablenotes}
\end{threeparttable}\end{table}


\paragraph{Deflactores.} Todas las regresiones salariales incluyen efectos fijos de trimestre, que absorben la inflación nacional, de modo que las estimaciones de diferencias en diferencias del logaritmo del salario no dependen del deflactor nacional; y las medidas de incidencia comparan salarios nominales con el piso legal nominal, por lo que no involucran deflactación alguna. Una deflactación de precios específica por departamento solo importaría para comparaciones entre departamentos de \emph{niveles} de salario real, que no cumplen ningún papel en la identificación.

\begin{figure}[!t]
\centering
\includegraphics[width=0.9\textwidth]{fig9_validation2025.pdf}
\caption{Estudio de eventos en torno a la reforma de enero de 2025 (validación fuera de muestra). Coeficientes de la interacción del indicador de baja calificación con indicadores de trimestre; trimestre de referencia 2024Q4. Las bandas son intervalos de confianza al noventa y cinco por ciento con errores agrupados por departamento.}
\label{fig:val2025}
\end{figure}

\begin{table}[!tbp]\centering
\caption{Reformas placebo en el tiempo (solo ventana previa a la reforma)}\label{tab:placebo}
\begin{threeparttable}
\begin{tabular}{lccc}
\toprule
Resultado & Placebo 2021Q2 & Placebo 2021Q3 & Placebo 2021Q4 \\
\midrule
Log salario real por hora & 0.010 & -0.026$^{*}$ & -0.023 \\
 & (0.025) & (0.013) & (0.015) \\
Empleo & -0.057$^{***}$ & -0.034$^{***}$ & -0.026$^{***}$ \\
 & (0.016) & (0.007) & (0.008) \\
Empleo formal & -0.024 & -0.029$^{*}$ & -0.017$^{**}$ \\
 & (0.015) & (0.014) & (0.007) \\
Trab. independiente & 0.027$^{*}$ & 0.002 & 0.019$^{*}$ \\
 & (0.014) & (0.009) & (0.009) \\
\bottomrule
\end{tabular}

\begin{tablenotes}\footnotesize
\item Notas: Cada celda es el coeficiente $\mathrm{Low}\times\mathrm{Post}$ de una estimación placebo de diferencias en diferencias que asigna una reforma ficticia al trimestre indicado, estimada usando solo datos previos a la reforma (2021Q1--2022Q1). Errores estándar agrupados por departamento entre paréntesis; todas las estimaciones se ponderan con los factores de expansión de la ENAHO. $^{*}p<0.1$, $^{**}p<0.05$, $^{***}p<0.01$.
\end{tablenotes}
\end{threeparttable}
\end{table}
